\chapter{Conclusions}
\label{chap:conclusion}

This thesis has developed and evaluated machine-learning methods to estimate the sensitivity of the fully hadronic vector boson scattering channel to anomalous quartic gauge couplings in the Éboli dimension-eight operator basis, on the full ATLAS Run-2 dataset of $140\,\mathrm{fb}^{-1}$ at $\sqrt{s}=13\,\mathrm{TeV}$. The fully hadronic channel carries the largest combined branching fraction of all VBS final states. The corresponding QCD multijet background exceeds the electroweak signal by roughly eight orders of magnitude in the analysis phase space, which makes a dedicated event-classification stage necessary before any sensitivity estimate is possible.

A two-stage classifier chain addresses this imbalance. The VBF-RNN tagger separates the VBS subprocess from the non-VBS electroweak $VVjj$ contributions to the same final state. Its six-jet version, which extends the input from the four tagging jets to include the two large-$R$ signal jets, raises the inclusive AUC from $0.760$ to $0.801$ over the four-jet version trained on the same target and removes the operator-family dependence that the earlier four-jet versions exhibited. At the working point $s>0.3$ used downstream, version~3 retains between $94\%$ and $98\%$ of true VBS events uniformly across the scalar, mixed and tensor operator families. An MLP discriminant then separates the aQGC signal from the composite background, taking twelve event-level features as input including the VBF-RNN score itself. Three architectures are evaluated: specialised binary classifiers trained per operator family, a global binary classifier covering all $19$ operators in the basis, and a four-node multi-class classifier resolving the scalar, mixed and tensor operator families. An equal-signal-cross-section reweighting addresses the operator-level imbalance inside each training category.

The Asimov significance of every discriminant on every probed operator is compared to the $m_{JJ}$-only baseline computed under the same binning prescription and $30\%$ Gaussian background uncertainty. The multivariate discriminants improve the expected significance by factors of $11$ to $281$ depending on the operator and architecture, with the largest relative gains on the rare scalar family where the $m_{JJ}$ baseline is most suppressed by the background uncertainty, and the smallest on the tensor family that already populates the high-$m_{JJ}$ tail. The equal-signal-cross-section reweighting is shown to be essential for the multi-class architecture: on the tensor family its significance relative to the specialised model rises from a range of $0.40$ to $0.53$ at the v3 configuration to a range of $0.79$ to $0.93$ at the v3 plus equal-$\sigma$ configuration adopted for the analysis. The dedicated working-point optimisation establishes that the chosen training selection of VBF-RNN${}>0.3$ and DisCoJet${}>0.65$ maximises the Asimov significance on the high-cross-section tensor and mixed operators that drive the analysis reach, at a modest and well-understood cost on the rare scalar operators.

The sensitivities reported include a conservative $30\%$ Gaussian background normalisation uncertainty and are derived from the pure-EFT quadratic-only signal model under the one-operator-at-a-time hypothesis. A definitive choice between the three classifier architectures requires profiling the remaining systematic uncertainties on the large-$R$ jet calibration, the subdominant backgrounds and the signal theory, and the headline limits require a unitarity-clipping prescription that bounds the dimension-eight expansion in the high-$m_{JJ}$ tail. The combined Run-2 plus Run-3 dataset of approximately $472\,\mathrm{fb}^{-1}$, which is now available, is expected to tighten the per-operator limits by approximately $26\%$ to $46\%$ depending on the operator. A subsequent combination of the present fully hadronic measurement with the existing semi-leptonic, fully-leptonic and triboson aQGC measurements provides a route to further tightening of the per-operator constraints below the current single-channel state of the art.
